% GROK REVIEW::GROK-2014::2026-10-14T20:14:00+09:00
\documentclass[11pt]{article}
\usepackage[margin=1in]{geometry}
\usepackage{amsmath,amssymb}
\title{Independent Grok Review --- 2026-10-14 20:14 KST}
\author{Grok}
\date{14 October 2026}

\begin{document}
\maketitle

\begin{center}
\fbox{\parbox{0.9\linewidth}{
\textbf{Executive status: P versus NP remains OPEN.}
No load-bearing implication is closed by this review. MERGEGUARD1 interfaces remain \CAND{} pending independent specialist review.
}}
\end{center}

\section{Scope}
This independent review examines the living document at repository-root \texttt{main.tex} (MERGEGUARD1 whole-document reconciliation, 2026-10-14), the MERGEGUARD1 synthesis snapshot and handoff, prior specialist reports, the guarded merging construction on the binary three-color kernel, and related PINCORE1 / BINARYPROJ1 / affine-residual materials. It confirms that current firewalls remain correctly classified and that the unrestricted whole-computation lower bound stays OPEN.

\section{Confirmation of current spine}
Affine base, residual contraction, cotree witness bijection, and upper-bound escapes remain PROVED/EXACT at their stated scopes. Prior cube selectors, TRACEISO1, and reviewed algebraic controls remain PROVED/EXACT. TWOTREEBRIDGE1, PINCORE1 pin-query closure, and BINARYPROJ1 Boolean-interior projection remain CANDIDATE. MERGEGUARD1's exact guarded quotient (polynomially many ternary exclusions yielding an exact image and full witness lift) and the polynomial 2-SAT check for prescribed simultaneous one/two-value fibers are retained as CANDIDATE pending independent review. Global unguarded merge-image safety is coNP-complete even on instances displaying a K4 contradiction; this is an extra certification task, not mandatory work for a source decider. Finding a successful fiber remains the open global choice problem. Growing-factor amplification remains CONDITIONAL. The unrestricted whole-computation lower bound remains OPEN.

No contradiction with prior exact results is identified. The guarded-merge construction eliminates certain spurious solutions without converting the source decision problem into a polynomial-time algorithm. Representation-driven structure continues to function as a useful firewall rather than a hardness source for unrestricted computation. Finite computational evidence is retained as COMPUTATIONAL and does not upgrade proof status.

\section{Open load-bearing arrows}
The implication from an arbitrary eligible short program (or uniform decider) to superpolynomial cost on a structured residual object remains OPEN. No model-matched all-algorithm theorem is established. Existing firewalls relativize and algebrize; no nonrelativizing hinge is known. Two-tree compatibility, pin-query closure, Boolean-interior projection, and guarded merging do not by themselves yield a superpolynomial lower bound for unrestricted computation. Opaque representations, arbitrary joint nonlinear compatibility, and historical compiler interfaces remain outside current selector theorems without additional bridges.

\section{Highest-value next target}
The single highest-value next assignment remains appropriate: implement a source-retained quotient search that learns an actual odd-cycle exclusion from a failed full block assignment, testing whether an independently justified contraction reduces unresolved choices while charging cycle discovery, learned-constraint storage, and original witness lifting. This avoids demanding an arbitrary coNP-complete merge-safety service. A qualitative contradiction suffices only if it excludes arbitrary polynomial deciders. Failure of one learning rule is not a lower bound for other algorithms. A successful selector or elimination is a useful firewall; it does not close a whole-machine implication.

\section{Status and recommendations}
P versus NP remains open. The living document has been updated to record this independent review. All claims continue to be classified under the allowed statuses. The comments/grok folder receives this review artifact. A dedicated comments folder structure for ongoing independent reviews is maintained under comments/grok.

% END_OF_VERIFIED_REVIEW::GROK::GROK-2014
\end{document}
